% !TeX root = ../papers/nc-p3.tex
\section{階数十の四次元中心に続く二次曲線と特異三項鉛筆}
\label{sec:rank-ten-conic-after-center}
\begin{theorem}\label{thm:rank-ten-conic-after-center}
反対階数10、曲面射影のない標準正則窓で、中央支持4の順中心 $p$ の第二曲線が
平面二次曲線であるとする。第一直線の唯一の順中心 $q$ は中央支持3であり、
その二項鉛筆は $L_1\oplus L_1^t$ 型である。第二曲線の次の第三曲線は平面に入り、
その次の位置も $p$ を通る平面に入る。第三曲線が三次なら平面三次の特異点は
$q$ の中央支持の消去点であり、そこに逆中心がある。
逆向きにも成立する。
\end{theorem}
\begin{proof}
前節と同じ $U=\ker p$ を用いる。第一直線と第二曲線の縮約を
$[s,t]A(t)$ とすると今度は $A=t_0A_0+t_1A_1$ が一次である。
第二位置の二次曲線の平面を法線 $h$ で表す。
定理\ref{thm:branch-rank-one-residual}の二次曲線の証明の通り、
第一直線上の三つの関係列を平面へ制限すると二項二次系の二つの標準一次関係と零列になる。
零列の元の値は $m(s)h$ で、$m$ は非零一次式である。
従って第一直線の唯一の順中心は $m=0$ の点 $q$ で、そこで他の二列は平面に制限しても独立である。
$q$ でのテンソルは、これら二つの独立な先頭文字に対応する $A_0,A_1$ の鉛筆に左右同値となる。
実際標準二列を第二曲線に評価すると $[s,t]t_0,[s,t]t_1$ であり、
列に対応する定数係数行列がそのまま $A_0,A_1$ である。
両行列を合わせた $U$ 支持は3（$q$ の前点は $p$ 一つ）、
位置3の支持次元を $k$ とするとこれは $q$ の中央支持である。

$D=A(t)V_3^*\subset U\otimes S_1$ は次元 $k$。
前節と同じ反対切断の制限から $\dim S_1D\le6$ である（今回は第二位置でも平面に制限する）。
もし $k=4$ なら $\dim\{u\in U:S_1u\subset D\}=8-\dim S_1D\ge2$ なので
$D=S_1U'$、$\dim U'=2$ となり $U$ 支持3に反する。
$k=2$ なら中央二項鉛筆の補題から、先頭と中央の二位置の三関係は二項対称二乗の平面である。
すなわち $A(t)$ の有効な三行二列を
$\left(\begin{smallmatrix}t_0&0\\ t_1&t_0\\0&t_1\end{smallmatrix}\right)$ とできる。
その左核行は $(t_1^2,-t_0t_1,t_0^2)$ で、非零の一次関係を持たない。
従って $([-,t]\otimes A(t)z)_{t,z}$ は六次元全体を張り、$R_0$ の第一直線への制限は単射となる。
しかし各 $\ker A(t)$ は二次元以上であり、単射ならその全点が $x_2(t)$ の後続となる。
曲線全体が順中心となって矛盾する。

従って $k=3$。一般の $t$ で $W_{x_2(t)}\subset W_{x_1(t)}^\perp$、
右辺を $W_p$ で割った次元が $10-4-4=2$ である
（第一直線の一般中央支持は前節と同じく4）。
$W_p$ は第一直線で零なので、別の $s$ で評価しても $A(t)$ は階数2以下となる。
有効な三行三列鉛筆は定数の左右核を持たず、従ってKronecker分解は
正の最小指数を持つ左右の特異ブロック各一つ、すなわち $L_1\oplus L_1^t$ である。
適合基底で
\[
 A(t)=u(t_0\alpha_0+t_1\alpha_1)^t+(t_0v_0+t_1v_1)b^t
\]
と書ける。$u,v_0,v_1$ は $U$ の基底、$\alpha_0,\alpha_1,b$ は位置3の独立な文字である。
$A(t)x_3(t)=0$ より $b(x_3)=0$ であり、$x_3$ は平面に入る。
$z$ を三文字の共通消去点とすると、$x_3$ は動くのでその射影は
$[\alpha_0(x_3):\alpha_1(x_3)]=[-t_1:t_0]$ である。
隣接する $x_3(t),x_4(t)$ を固定して全ての $y$ で
$x_4(t)^t A(y)x_3(t)=0$ を使うと $u(x_4(t))=0$ となる。
平面は $p$ を含む。

第三曲線が三次なら $z$ からの射影次数1なので $z$ は重複度2の点である。
そこで逆評価が階数3なら、逆辺の三次小行列の非零な一つを使い、近傍で前点が
$x_3$ の正則関数として定まる。正規化パラメータでの合成は滑らかな二次曲線
$x_2$ への同型である。従って $z$ の二つの逆像は同じ点へ写せず、
一つの逆像であっても正規化の微分が零ではあり得ない。
平面三次の二重点では二枝または一枝で接写像が零なので矛盾する。
逆評価はちょうど階数2である。
\end{proof}

\begin{proposition}\label{prop:singular-ternary-center-projections}
反対階数によらず、中央支持3の順中心 $q$ が $L_1\oplus L_1^t$ 型なら、
その次の位置のさらに次の位置と、その後の位置における二つの平面射影が
一点ずつからの射影として同定される。ただし共通の消去平面に入らない点対での主張である。
特に $q$ の列の第二曲線がねじれ三次曲線なら第三次数は2以上4以下である。
その三次曲線が $q$ の中央支持の消去点を通るときは上界が3になる。
\end{proposition}
\begin{proof}
$q$ の位置を1として上の鉛筆を読む。ここで $u,v_0,v_1$ は $q$ の末尾の三支持であり、
その全消去点は $q$ の一意な前点である。位置3,4の辺の二関係は
$u(x_4)\alpha_j(x_3)+v_j(x_4)b(x_3)=0$ ($j=0,1$)。
従って $b(x_3)\ne0$ で $x_4$ が全消去点でなければ
\[
 [u(x_4):v_0(x_4):v_1(x_4)]
       =[b(x_3):-\alpha_0(x_3):-\alpha_1(x_3)].
\]
$q$ 自身の第二曲線が全空間を張れば右の像は直線に入らない。
第三曲線が動けば左の射影も一般に定義され、元の曲線は直線ではあり得ない。
第二曲線がねじれ三次ならその直後に一定の逆中心へ入ることもない。
このとき右辺の射の引戻し次数は、射影の中心が三次曲線上にあれば2、なければ3である。
左辺の射影の中心 $q^-$ は $q$ の一意な前点で、その第三曲線による逆像長は高々1である。
実際相異なる二点なら、それぞれ $q$ から第三点 $q^-$ までの経路に辺 $q^-\to q$ を加え、
独立に反復できる二つの閉経路ができる。一点に長さ2があれば同じ経路の双対数族があり、
最初の直線で非零の微分を持つ。各反復で独立な双対数を取ると、各ブロックの一次文字の値が
その二次元環を張るので、長さ $4n$ の経路の評価次元が $2^n$ 以上になる。
いずれも標準成長に反する（命題\ref{prop:terminal-branch-contact}の帰路の議論）。
従って左辺の射の引戻し次数に足す基点因子の長さは高々1であり、表示した上下界が従う。
\end{proof}
